\subsection{Finite-feature transfer of negative curvature at a fixed time}
\label{subsection:concentration-kappa}

This fixed-time supplement is not an input to the tracking theorem,
objective comparison, or fixed-time gradient comparison.
Its finite-feature conclusion does not control whole-path curvature-response
factors or adaptive jump times.

We give a sufficient sample size for a negative population Hessian
direction to remain negative in the reference sample. The statement
uses an exact finite-feature representation and a spatial-Hessian
condition. We first fix a deterministic law; the conditional
interpretation for a fixed time of PWGF follows after the proof.

\paragraph{Finite features and the two Hessian terms.}
Consider the finite-feature case of
\cref{assumption-structured-moment-smoothness}: for admissible laws $\nu$,
\begin{equation}
 F(\nu)=G\left(\int\phi\,d\nu\right),\qquad
 \phi:\mathbb R^d\longrightarrow\mathbb R^q,\qquad q\in\mathbb N.
 \label{eq:finite-feature-objective}
\end{equation}
The dimension $q$ is fixed independently of $N$; no feature cutoff or
approximation of $F$ is taken. At the law $\mu$ under consideration,
define
\begin{equation}
 \Phi(x):=\fD\phi(x)\in\mathbb R^{q\times d},\qquad
 A:=\fD^2G\left(\int\phi\,d\mu\right)=A^\top\in\mathbb R^{q\times q}.
 \label{eq:finite-feature-derivatives}
\end{equation}
The moment chain rule in \cref{eq:weighted-integral-hessian-representation}
reduces to the finite sum
\[
 \nabla_W^2F(\mu;x,y)
 =\sum_{j,k=1}^q A_{jk}\nabla\phi_j(x)\nabla\phi_k(y)^\top
 =\Phi(x)^\top A\Phi(y).
\]
Recall the integral and spatial Hessian operators from
\cref{subsec:population-pwgf}. On $L^2(\mu;\mathbb R^d)$ they are
\begin{equation}
 \begin{split}
 (H_\mu v)(x)&:=
   \int\Phi(x)^\top A\Phi(y)v(y)\,\mu(dy),\\
 (H'_\mu v)(x)&:=
   \nabla_x\nabla_WF(\mu,x)v(x).
 \end{split}
 \label{eq:finite-feature-hessian-operators}
\end{equation}
Their sum is the full operator $\mathcal H_\mu=H_\mu+H'_\mu$.
For independent $Y_1,\ldots,Y_N\sim\mu$, its sampled full Hessian is
\begin{equation}
 \begin{split}
 \mathcal H^{Y,N}(i,j)
 &:=\frac1N\Phi(Y_i)^\top A\Phi(Y_j)
      +\delta_{ij}\nabla_x\nabla_WF(\mu,Y_i),\\
 \kappa^{Y,N}&:=\lambda_{\min}(\mathcal H^{Y,N}).
 \end{split}
 \label{eq:finite-feature-sampled-hessian}
\end{equation}
The matrix acts on $(\mathbb R^d)^N$ with its Euclidean inner product.
At $\mu=\mu_t$, this is precisely the reference-sample curvature
$\kappa_t^{Y,N}$ of \eqref{eq:curvature-main}: derivatives are evaluated
at $\mu_t$, not at the empirical measure.

\Needspace{8\baselineskip}
\begin{proposition}[Transfer of negative curvature to the reference sample]
\label{prop:finite-feature-negative-curvature}
Fix $\mu$ as above, and suppose that the following hold for deterministic
constants $K\ge1$, $C_s\ge0$, $\varepsilon\ge0$, and $\gamma>0$.
\begin{enumerate}
 \item The coefficient and feature moments satisfy
 \begin{equation}
 \|A\|_{\mathbb R^q\to\mathbb R^q}\le K,\qquad
 \mathbb E\|\Phi(Y_1)\|_F^4\le K^4.
 \label{eq:finite-feature-common-bound}
 \end{equation}
 Here $\|\cdot\|_F$ is the Frobenius norm, the square root of the sum of
 squared entries.
 \item The spatial Hessian is controlled by the population gradient:
 \begin{equation}
 \mathop{\rm ess\,sup}_{x\sim\mu}
 \|\nabla_x\nabla_WF(\mu,x)\|_{\mathbb R^d\to\mathbb R^d}
 \le C_s\|\nabla_WF(\mu,\cdot)\|_{L^2(\mu;\mathbb R^d)}.
 \label{eq:finite-feature-spatial-condition}
 \end{equation}
 \item Near-stationarity and negative full curvature hold separately:
 \begin{equation}
 \begin{split}
 \|\nabla_WF(\mu,\cdot)\|_{L^2(\mu;\mathbb R^d)}&\le\varepsilon,
       \qquad C_s\varepsilon\le\gamma/8,\\
 \inf\operatorname{Spec}(H_\mu+H'_\mu)&\le-\gamma.
 \end{split}
 \label{eq:finite-feature-negative-gap}
 \end{equation}
\end{enumerate}
Then, for every $\delta\in(0,1)$,
\begin{equation}
 N\ge\frac1\delta\max\left\{1,\frac{12^4K^{12}}{\gamma^4}\right\}
 \quad\Longrightarrow\quad
 \mathbb P(\kappa^{Y,N}\le-\gamma/2)\ge1-\delta.
 \label{eq:finite-feature-negative-transfer}
\end{equation}
\end{proposition}

Condition \eqref{eq:finite-feature-spatial-condition} has the form of
Assumption~2, equation~(8), in \citet{yamamoto2025hessianguided}.
That work imposes the condition for every law in its setting; here it is
required only at the law in the proposition. It does not follow merely
from finite features or from a small gradient.

\begin{proof}
\noindent\textbf{Step 1: a single covariance event.}
Define the population and sample feature covariances
\begin{equation}
 \Sigma:=\mathbb E[\Phi(Y_1)\Phi(Y_1)^\top],\qquad
 \widehat\Sigma:=\frac1N\sum_{i=1}^N\Phi(Y_i)\Phi(Y_i)^\top.
 \label{eq:finite-feature-covariances}
\end{equation}
They are positive semidefinite $q\times q$ matrices. If $s_j$ are the
singular values of a matrix $\Phi$, then
\[
 \|\Phi\Phi^\top\|_F^2=\sum_js_j^4
 \le\left(\sum_js_j^2\right)^2=\|\Phi\|_F^4.
\]
The centered matrices
$Z_i:=\Phi(Y_i)\Phi(Y_i)^\top-\Sigma$ are independent and have mean zero.
Expanding the squared Frobenius norm gives
\[
 \mathbb E\|Z_i\|_F^2
 =\mathbb E\|\Phi(Y_i)\Phi(Y_i)^\top\|_F^2-\|\Sigma\|_F^2
 \le K^4.
\]
For $i\ne j$, independence gives
$\mathbb E\langle Z_i,Z_j\rangle_F
=\langle\mathbb EZ_i,\mathbb EZ_j\rangle_F=0$.
Consequently
\begin{align*}
 \mathbb E\|\widehat\Sigma-\Sigma\|_F^2
 &=\frac1{N^2}\sum_{i,j}\mathbb E\langle Z_i,Z_j\rangle_F\\
 &=\frac1{N^2}\sum_i\mathbb E\|Z_i\|_F^2\le\frac{K^4}{N}.
\end{align*}
Markov's inequality, applied to this squared norm, yields
\[
 \mathbb P\left(\|\widehat\Sigma-\Sigma\|_F>
                   \frac{K^2}{\sqrt{N\delta}}\right)
 \le\frac{K^4/N}{K^4/(N\delta)}=\delta.
\]
Thus, with probability at least $1-\delta$,
\begin{equation}
\label{eq:finite-feature-covariance-success}
 \|\widehat\Sigma-\Sigma\|_{\mathbb R^q\to\mathbb R^q}
 \le\|\widehat\Sigma-\Sigma\|_F
 \le\frac{K^2}{\sqrt{N\delta}}.
\end{equation}
All subsequent sample calculations use
\eqref{eq:finite-feature-covariance-success}.
Since $\|\Phi\Phi^\top\|_{\mathbb R^q\to\mathbb R^q}\le\|\Phi\|_F^2$,
\[
 \|\Sigma\|_{\mathbb R^q\to\mathbb R^q}\le\mathbb E\|\Phi(Y_1)\|_F^2
 \le(\mathbb E\|\Phi(Y_1)\|_F^4)^{1/2}\le K^2.
\]
The sample-size condition gives $N\delta\ge1$, and hence, on
\eqref{eq:finite-feature-covariance-success},
\[
 \|\widehat\Sigma\|_{\mathbb R^q\to\mathbb R^q}
 \le\|\Sigma\|_{\mathbb R^q\to\mathbb R^q}+\|\widehat\Sigma-\Sigma\|_{\mathbb R^q\to\mathbb R^q}
 \le K^2+\frac{K^2}{\sqrt{N\delta}}\le2K^2.
\]

\par\noindent\textbf{Step 2: compare the square roots and Hessian matrices.}
Continue on the covariance bound \eqref{eq:finite-feature-covariance-success}.
Take a unit eigenvector $v$ of the symmetric matrix
$\widehat\Sigma^{1/2}-\Sigma^{1/2}$ with eigenvalue $\lambda$ such that
$|\lambda|=\|\widehat\Sigma^{1/2}-\Sigma^{1/2}\|_{\mathbb R^q\to\mathbb R^q}$.
Here each square root is the positive semidefinite square root.
The identity
\[
 \widehat\Sigma-\Sigma
 =\widehat\Sigma^{1/2}(\widehat\Sigma^{1/2}-\Sigma^{1/2})
  +(\widehat\Sigma^{1/2}-\Sigma^{1/2})\Sigma^{1/2}
\]
holds even when the two matrices do not commute.
Positive semidefiniteness gives
\[
 v^\top(\widehat\Sigma^{1/2}+\Sigma^{1/2})v
 \ge|v^\top(\widehat\Sigma^{1/2}-\Sigma^{1/2})v|=|\lambda|.
\]
Multiplying the preceding identity by $v^\top$ and $v$, and using the
eigenvector equation on its two sides, gives
\begin{align*}
 \|\widehat\Sigma-\Sigma\|_{\mathbb R^q\to\mathbb R^q}
 &\ge|v^\top(\widehat\Sigma-\Sigma)v|\\
 &=|\lambda|\,v^\top(\widehat\Sigma^{1/2}+\Sigma^{1/2})v\\
 &\ge|\lambda|^2
   =\|\widehat\Sigma^{1/2}-\Sigma^{1/2}\|_{\mathbb R^q\to\mathbb R^q}^2.
\end{align*}
Thus Step 1 implies
$\|\widehat\Sigma^{1/2}-\Sigma^{1/2}\|_{\mathbb R^q\to\mathbb R^q}
\le K(N\delta)^{-1/4}$.
Define the symmetric matrices reused in the eigenvalue comparison by
\[
 \mathcal T:=\Sigma^{1/2}A\Sigma^{1/2},\qquad
 \widehat{\mathcal T}:=\widehat\Sigma^{1/2}A\widehat\Sigma^{1/2}.
\]
Inserting and subtracting $\Sigma^{1/2}A\widehat\Sigma^{1/2}$ gives
\begin{align*}
 \widehat{\mathcal T}-\mathcal T
 &=(\widehat\Sigma^{1/2}-\Sigma^{1/2})A\widehat\Sigma^{1/2}
       +\Sigma^{1/2}A(\widehat\Sigma^{1/2}-\Sigma^{1/2}),\\
 \|\widehat{\mathcal T}-\mathcal T\|_{\mathbb R^q\to\mathbb R^q}
 &\le K(N\delta)^{-1/4}\,K\,(\sqrt2K+K)\\
 &\le3K^3(N\delta)^{-1/4}\le\gamma/4.
\end{align*}
The last inequality follows directly from
$(N\delta)^{1/4}\ge12K^3/\gamma$ in
\eqref{eq:finite-feature-negative-transfer}.

\par\noindent\textbf{Step 3: identify the negative eigenvalues.}
Define the bounded feature operator and its adjoint by
\[
 \begin{aligned}
 \mathcal B &:L^2(\mu;\mathbb R^d)\to\mathbb R^q,
 &\mathcal Bv&=\int\Phi(x)v(x)\,\mu(dx),\\
 \mathcal B^* &:\mathbb R^q\to L^2(\mu;\mathbb R^d),
 &(\mathcal B^*u)(x)&=\Phi(x)^\top u.
 \end{aligned}
\]
Indeed, Cauchy--Schwarz gives
\begin{align*}
 \|\mathcal Bv\|_{\mathbb R^q}
 &\le\int\|\Phi(x)\|_F\|v(x)\|_{\mathbb R^d}\,\mu(dx)\\
 &\le(\mathbb E\|\Phi(Y_1)\|_F^2)^{1/2}\|v\|_{L^2(\mu;\mathbb R^d)}
 \le K\|v\|_{L^2(\mu;\mathbb R^d)}.
\end{align*}
Moreover
$\langle\mathcal Bv,u\rangle_{\mathbb R^q}
=\int\langle v(x),\Phi(x)^\top u\rangle_{\mathbb R^d}\,\mu(dx)$,
which verifies the adjoint formula. By
\cref{eq:finite-feature-covariances,eq:finite-feature-hessian-operators},
for $u\in\mathbb R^q$ and $v\in L^2(\mu;\mathbb R^d)$, substitution gives
\begin{align*}
 \mathcal B\mathcal B^*u
 &=\int\Phi(x)\Phi(x)^\top u\,\mu(dx)=\Sigma u,\\
 (\mathcal B^*A\mathcal Bv)(x)
 &=\Phi(x)^\top A\int\Phi(y)v(y)\,\mu(dy)=(H_\mu v)(x).
\end{align*}
Thus $\mathcal B\mathcal B^*=\Sigma$ and
$H_\mu=\mathcal B^*A\mathcal B$.
In particular $H_\mu$ is self-adjoint and has finite rank.

We use the following elementary eigenvector calculation for products of
two bounded maps $U,V$ between two spaces. For $\lambda\ne0$,
\begin{align*}
 UVx=\lambda x,\quad x\ne0
 &\ \Longrightarrow\
 Vx\ne0,\quad VU(Vx)=\lambda Vx,\\
 VUy=\lambda y,\quad y\ne0
 &\ \Longrightarrow\
 Uy\ne0,\quad UV(Uy)=\lambda Uy.
\end{align*}
For the first implication, $Vx=0$ would give $\lambda x=UVx=0$,
a contradiction; the second is proved in the same way.
Apply these identities first with $U=\mathcal B^*$, $V=A\mathcal B$,
then with $U=A\Sigma^{1/2}$, $V=\Sigma^{1/2}$:
\[
 \mathcal B^*A\mathcal B
 \ \longleftrightarrow\ A\Sigma
 \ \longleftrightarrow\ \Sigma^{1/2}A\Sigma^{1/2}=\mathcal T
 \quad\hbox{for nonzero eigenvalues}.
\]
To check why only eigenvalues matter here, take
$v\perp\operatorname{ran}H_\mu$. Self-adjointness implies, for every $w$,
\[
 \langle H_\mu v,w\rangle_{L^2(\mu;\mathbb R^d)}
 =\langle v,H_\mu w\rangle_{L^2(\mu;\mathbb R^d)}=0,
 \qquad\text{hence }H_\mu v=0.
\]
The range is finite-dimensional and thus closed, so
\[
 L^2(\mu;\mathbb R^d)
 =\operatorname{ran}H_\mu\oplus(\operatorname{ran}H_\mu)^\perp,
 \qquad
 H_\mu=(H_\mu|_{\operatorname{ran}H_\mu})\oplus0.
\]
The first restriction is a finite-dimensional self-adjoint matrix.
Consequently any negative spectral bottom is one of its eigenvalues.
This justifies
using the displayed eigenvector correspondence for its spectral bottom
whenever that bottom is negative.

For the sample, set
$P=N^{-1/2}[\Phi(Y_1)\ \cdots\ \Phi(Y_N)]\in\mathbb R^{q\times dN}$.
Matrix multiplication gives
\[
 PP^\top=\widehat\Sigma,\qquad
 (P^\top AP)(i,j)=N^{-1}\Phi(Y_i)^\top A\Phi(Y_j).
\]
The same two product calculations, replacing $\mathcal B$ by $P$, show
that the negative eigenvalues of $P^\top AP$ and
$\widehat{\mathcal T}$ coincide. No assertion about the multiplicity
of zero eigenvalues is needed.

\par\noindent\textbf{Step 4: restore both spatial-Hessian terms.}
By \cref{eq:finite-feature-spatial-condition,eq:finite-feature-negative-gap},
\begin{gather*}
 \|H'_\mu v\|_{L^2(\mu;\mathbb R^d)}^2
 \le(C_s\varepsilon)^2\|v\|_{L^2(\mu;\mathbb R^d)}^2,\\
 \|H'_\mu\|_{L^2(\mu;\mathbb R^d)\to L^2(\mu;\mathbb R^d)}\le C_s\varepsilon\le\gamma/8.
\end{gather*}
For every unit $v\in L^2(\mu;\mathbb R^d)$, subtracting the spatial term gives
\begin{align*}
 \langle v,H_\mu v\rangle_{L^2(\mu;\mathbb R^d)}
 &=\langle v,(H_\mu+H'_\mu)v\rangle_{L^2(\mu;\mathbb R^d)}
      -\langle v,H'_\mu v\rangle_{L^2(\mu;\mathbb R^d)}\\
 &\le\langle v,(H_\mu+H'_\mu)v\rangle_{L^2(\mu;\mathbb R^d)}+\gamma/8.
\end{align*}
Taking infima of the Rayleigh quotients yields
\[
 \inf\operatorname{Spec}(H_\mu)
 \le\inf\operatorname{Spec}(H_\mu+H'_\mu)+\gamma/8
 \le-7\gamma/8.
\]
The negative spectral bottom of $H_\mu$ is an eigenvalue.
Step 3 therefore gives $\lambda_{\min}(\mathcal T)\le-7\gamma/8$.
For a unit minimum-eigenvalue vector $u$ of $\mathcal T$,
\[
 \lambda_{\min}(\widehat{\mathcal T})
 \le u^\top\widehat{\mathcal T}u
 \le u^\top\mathcal Tu+\|\widehat{\mathcal T}-\mathcal T\|_{\mathbb R^q\to\mathbb R^q}
 \le-7\gamma/8+\gamma/4=-5\gamma/8.
\]
Its negative eigenvalue corresponds to one of $P^\top AP$ by Step 3.

The essential bound in \eqref{eq:finite-feature-spatial-condition}
holds at every $Y_i$ almost surely, since a finite union of $\mu$-null
sets still has sampling probability zero. The block-diagonal spatial
matrix has norm at most
$\max_i\|\nabla_x\nabla_WF(\mu,Y_i)\|_{\mathbb R^d\to\mathbb R^d}\le\gamma/8$.
Evaluating the sampled full Hessian on a minimum-eigenvalue vector of
$P^\top AP$ now gives
\[
 \kappa^{Y,N}
 \le\lambda_{\min}(P^\top AP)
       +\max_i\|\nabla_x\nabla_WF(\mu,Y_i)\|_{\mathbb R^d\to\mathbb R^d}
 \le-5\gamma/8+\gamma/8=-\gamma/2.
\]
The only event with nonzero failure probability is the failure of
\eqref{eq:finite-feature-covariance-success}, with probability at most $\delta$.
\end{proof}

\paragraph{Fixed-time conditional interpretation.}
At a deterministic time $t$, condition on the population fields
$\mathcal G_{\rm pop}$ from \cref{lemma-gp-jump-implies-conditional-iid}
and put $\mu=\mu_t$. Let $E\in\mathcal G_{\rm pop}$ be any event on
which
\cref{eq:finite-feature-common-bound,eq:finite-feature-spatial-condition,eq:finite-feature-negative-gap}
hold with the same deterministic $q,K,C_s,\varepsilon,\gamma$.
Here $\Phi=\fD\phi$ is measure-independent, while
$A=\fD^2G(\int\phi\,d\mu_t)$ and all Hessian terms are evaluated at
the population law $\mu_t$.
For $N$ as in \eqref{eq:finite-feature-negative-transfer}, the
conditional-i.i.d. property and the tower rule give
\[
 \mathbb P(E\cap\{\kappa_t^{Y,N}>-\gamma/2\})
 =\mathbb E\!\left[\mathbf1_E
    \mathbb P(\kappa_t^{Y,N}>-\gamma/2\mid\mathcal G_{\rm pop})\right]
 \le\delta\,\mathbb P(E)\le\delta.
\]
This statement does not estimate $\mathbb P(E)$ or a whole-path curvature-response factor,
and it concerns the reference sample $Y$, not the empirical dynamics $X$.

\paragraph{A simple instance of the assumptions.}
In \eqref{eq:canonical-objective}, take $a=0$ and a law of mean zero.
The one-feature representation is $\phi(x)=x$ and
$G(z)=\beta\cos(\pi z)$. Thus $\Phi=1$, $A=-\beta\pi^2$,
$\nabla_WF(\mu,x)=0$, and $\nabla_x\nabla_WF(\mu,x)=0$.
The constant unit function is an eigenvector of $H_\mu$ with
eigenvalue $-\beta\pi^2$. Hence the conditions hold with
$K=\max\{1,\beta\pi^2\}$, $C_s=\varepsilon=0$, and
$\gamma=\beta\pi^2$. This observation does not impose the additional
spatial-Hessian condition on the variance-plus-cosine example with general $a$.
